\section*{الفصل \ref{ch:g10:habitable-planet} --- كوكب صالح للحياة}

\begin{solution}{exo:g10:habitable-planet:1}
هي درجة الحرارة التي يشع عندها باعث مثالي نحو الفضاء القدرة نفسها تمامًا التي
يمتصها الكوكب من نجمه. وسطح الأرض أدفأ
(\qty{288}{K} مقابل نحو \qty{255}{K}) لأن غازات الدفيئة تمتص جزءًا من الأشعة
تحت الحمراء التي يطلقها وتعيد جزءًا منها نحو الأسفل.
\end{solution}

\begin{solution}{exo:g10:habitable-planet:2}
الزهرة: \qty{737}{K} مقيسة مقابل \omterm{def:g10:habitable-planet:equilibrium-temperature}{درجة حرارة توازن} تقارب
\qty{227}{K}، أي تدفئة بنحو \qty{510}{K}.
\end{solution}

\begin{solution}{exo:g10:habitable-planet:3}
$S = 1361/1.52^2 \approx \qty{589}{W/m^2}$ (وتعطي بطاقة البيانات
\qty{586}{W/m^2}، بالمسافة المتوسطة الدقيقة).
\end{solution}

\begin{solution}{exo:g10:habitable-planet:4}
\qty{1}{bar}، \qty{330}{K}: سائل (دون نقطة الغليان، نحو
\qty{373}{K}). \qty{0.001}{bar}، \qty{250}{K}: جليد (فالضغط يقع فوق منحنى
الجليد--البخار، الذي يبلغ نحو \qty{0.0008}{bar} عند \qty{250}{K}).
\qty{100}{bar}، \qty{700}{K}: بخار --- فدرجة الحرارة أعلى من درجة الحرارة
الحرجة، فلا يمكن أن يوجد سائل.
\end{solution}

\begin{solution}{exo:g10:habitable-planet:5}
كوكب يدور حول نجم غير الشمس. العبور: يمر الكوكب أمام نجمه فيخفت ضوؤه قليلًا في
كل دورة. والتمايل: تحرّك جاذبيةُ الكوكب النجمَ ذهابًا وإيابًا على امتداد خط
نظرنا، فيُزاح لون ضوئه دوريًّا.
\end{solution}

\begin{solution}{exo:g10:habitable-planet:6}
$T = (1361 \times 0.706/(4 \times 5.67\times10^{-8}))^{1/4} \approx
\qty{255}{K} \approx \qty{-18}{\degreeCelsius}$. ومع $A = 0$:
$\approx \qty{278}{K} \approx \qty{5}{\degreeCelsius}$.
\end{solution}

\begin{solution}{exo:g10:habitable-planet:7}
$\sqrt 4 = 2$: من $2 \times 0.99 = \qty{1.98}{au}$ إلى $2 \times 1.70 =
\qty{3.40}{au}$. وعلى بعد \qty{1}{au} يتلقى الكوكب أربعة أضعاف إشعاعية
الأرض: فهو دون الحافة الداخلية بكثير، وحار جدًا.
\end{solution}

\begin{solution}{exo:g10:habitable-planet:8}
$S = 1361/0.723^2 \approx \qty{2604}{W/m^2}$؛
$T_{\mathrm{eq}} = (2604 \times 0.23/(4\sigma))^{1/4} \approx \qty{227}{K}$.
والسطح أسخن بنحو \qty{510}{K}: فغلافها الجوي الكثيف من ثاني أكسيد الكربون
(\qty{92}{bar}) ينتج أثر دفيئة هائلًا.
\end{solution}

\begin{solution}{exo:g10:habitable-planet:9}
عند \qty{6.36}{mbar} لا يكون الماء سائلًا إلا بين نحو \qty{0}{\degreeCelsius}
و\qty{0.5}{\degreeCelsius}. وعند \qty{10}{\degreeCelsius} يكون الضغط دون منحنى
الغليان: فيغلي الماء ويتبخر؛ وفي الليل يتجمد.
\end{solution}

\begin{solution}{exo:g10:habitable-planet:10}
$11.2/2.4 \approx 4.7$. فعلى جزيء الغاز أن يتحرك أسرع بخمس مرات تقريبًا كي
يفلت من الأرض مقارنةً بالقمر؛ وعلى مدى مليارات السنين تركت جاذبية القمر
الضعيفة غازاته تفلت، بينما احتفظت جاذبية الأرض بمعظم غازاتها.
\end{solution}

\begin{solution}{exo:g10:habitable-planet:11}
$T = (1361 \times 0.50/(4\sigma))^{1/4} \approx \qty{234}{K}$، أي نحو
\qty{-39}{\degreeCelsius}. وسيعكس الجليد من ضوء الشمس قدرًا كبيرًا يبقى معه
الكوكب باردًا بما يكفي للاحتفاظ به: إنها حلقة تعزّز نفسها.
\end{solution}

\begin{solution}{exo:g10:habitable-planet:12}
$T_{\mathrm{eq}} = (S_0(1-A)/(4\sigma d^2))^{1/4} = (S_0/(4\sigma))^{1/4}
(1-A)^{1/4} (d^2)^{-1/4}$؛ و$(1361/(4\sigma))^{1/4} \approx \qty{278.3}{K}$
و$(d^2)^{1/4} = \sqrt d$. وعلى بعد أربعة أضعاف: تُقسم $T$ على $\sqrt 4 = 2$.
\end{solution}

\begin{solution}{exo:g10:habitable-planet:13}
$\sqrt d = 278.3 \times 0.706^{1/4}/273 \approx 0.934$، إذن $d \approx
\qty{0.873}{au}$، أي دون الحافة الداخلية (\qty{0.99}{au}). \omterm{def:g10:habitable-planet:equilibrium-temperature}{فدرجة حرارة التوازن}
تهمل أثر الدفيئة: والكوكب الحقيقي ذو الغلاف الجوي أدفأ من $T_{\mathrm{eq}}$،
فيسخن أكثر مما ينبغي ويفقد ماءه على بعد أكبر، نحو \qty{0.99}{au}. فالحافتان
تأتيان من نماذج للغلاف الجوي كله، لا من $T_{\mathrm{eq}}$ وحدها.
\end{solution}

\begin{solution}{exo:g10:habitable-planet:14}
بين نقطة الانصهار، نحو \qty{0}{\degreeCelsius}، ونقطة الغليان، نحو
\qty{7}{\degreeCelsius}. وعند \qty{15}{\degreeCelsius} تغلي
البحيرة.
\end{solution}

\begin{solution}{exo:g10:habitable-planet:15}
$\sqrt{0.5} \approx 0.707$: من $0.707 \times 0.99 \approx \qty{0.70}{au}$
إلى $0.707 \times 1.70 \approx \qty{1.20}{au}$. فالكوكب على الحافة الخارجية
تمامًا. والحافتان نتيجتا نماذج والتغيير بالمقياس تقريبي، فيتعذر القول هل هو
داخل النطاق أم خارجه.
\end{solution}

\begin{solution}{pb:g10:habitable-planet:1}
\textbf{1.} $0.00151/0.0485^2 \approx 0.642$.

\textbf{2.} $0.642 \times 1361 \approx \qty{874}{W/m^2}$.

\textbf{3.} بين الاثنتين: أقل من إشعاعية الزهرة بنحو ثلاث مرات، وأكثر من إشعاعية
المريخ بمرة ونصف.

\textbf{4.} $365.25/11.18 \approx 32.7$.

\textbf{5.} $(874/(4 \times 5.67\times10^{-8}))^{1/4} \approx \qty{249}{K}$.

\textbf{6.} $(874 \times 0.23/(4\sigma))^{1/4} \approx \qty{173}{K}$.

\textbf{7.} ولا واحدة منهما: فالقيمة \qty{249}{K} ما زالت دون \qty{273}{K}. والغطاء
اللامع من السحب أو الجليد يخفض \omterm{def:g10:habitable-planet:equilibrium-temperature}{درجة حرارة التوازن} بأكثر من
\qty{75}{K}: فما يغطي الكوكب لا يقل أهمية عن بعده.

\textbf{8.} $\sqrt{0.00151} \approx 0.0389$.

\textbf{9.} $0.99 \times 0.0389 \approx \qty{0.0385}{au}$.

\textbf{10.} $1.70 \times 0.0389 \approx \qty{0.066}{au}$.

\textbf{11.} نعم: $0.0385 < 0.0485 < 0.066$.

\textbf{12.} $d_\odot = 0.0485/0.0389 \approx \qty{1.25}{au}$: في النصف
الداخلي من \omterm{def:g10:habitable-planet:habitable-zone}{النطاق الصالح للحياة} حول الشمس، بين الأرض والمريخ.

\textbf{13.} $1.055^{1/3} \approx 1.018$.

\textbf{14.} $11.2 \times \sqrt{1.055/1.018} \approx \qty{11.4}{km/s}$، أي قريبًا
من سرعة الإفلات من الأرض: فجاذبيته قادرة على الاحتفاظ بغلاف جوي (أما هل له غلاف
فعلًا فيتوقف على تاريخه).

\textbf{15.} لا يكشف تمايل النجم إلا جزء حركته الواقع على امتداد خط نظرنا؛ وبما
أن ميل المدار مجهول، فالكتلة المستنتجة لا تساوي الكتلة الحقيقية إلا إذا رُئي
المدار من حافته، وهي حدّ أدنى في غير ذلك.

\textbf{16.} ضغط النقطة الثلاثية، \qty{611.657}{Pa}، أي نحو
\qty{0.006}{bar}.

\textbf{17.} $(874 \times 0.706/(4\sigma))^{1/4} \approx \qty{228}{K}$.

\textbf{18.} $288 - 228 = \qty{60}{K}$: يحتاج بروكسيما~ب إلى تدفئة بأثر الدفيئة
قدرها \textbf{نحو \qty{60}{K}}، أي قرابة ضعف القيمة \qty{33}{K} على الأرض.
\end{solution}
